\section{Closing a Financing Gap}\label{sec:financing}

Population aging puts the financing of Social Security under pressure \citep{trustees2024,cbo2024socialsecurity}, and closing a financing gap requires raising revenue or cutting outlays. The MVPF ranks the instruments available for this purpose by the welfare cost of each dollar they save \citep{hendren2016policy,hendren2020unified}. This section draws the implications of our estimates for that choice; all statements refer to the reforms in our grid and to the cohort we study.

\paragraph{Reforms that save money and raise welfare.} A reform with an infinite MVPF relaxes the budget at no welfare cost and therefore comes first in any consolidation. In our grid, these are the reforms that reward later claiming (\autoref{tab:claiming}). Their savings are not small. The DRC increases lower the present value of government spending by 0.44 to 0.78 percent, and the symmetric adjustments by 0.12 percent (AL1) to 1.57 percent (AL9), while raising welfare by 0.10 to 0.80 percent. Within the symmetric family, fiscal saving and welfare rise together with the slope of the schedule, so the reform that saves the most, AL9, also raises welfare the most. For these reforms there is no trade-off between the budget and welfare to resolve: the government obtains the saving by paying a larger reward to those who delay, and the induced delays and longer careers more than finance the reward (Section~\ref{sec:results_claiming}).

\paragraph{Rewarding delay and raising the NRA.} The relevant comparison is with increases in the NRA, the most common response to financing pressure. For a fiscal saving up to the size of AN1, 0.87 percent of government spending, several reward-delay reforms deliver a similar or larger saving with a welfare gain instead of a welfare loss. AD5 saves 0.78 percent and AL4 saves 0.86 percent; AL5, AL6, AL7, AL8, and AL9 save 1.17, 1.37, 1.38, 1.47, and 1.57 percent. AN1 lowers welfare by 0.38 percent; each of these reforms raises it. Larger consolidations are a different matter. AN2 and AN3 save 1.75 and 2.38 percent, more than any single reward-delay reform in our grid. We have not computed combinations of reforms, such as a steeper claiming schedule together with a higher NRA, so the grid does not say how a saving of that size would best be obtained. It does show that savings up to 1.57 percent of government spending can be obtained while raising welfare, and that, for the cohort we study, the steepest reward-delay reform saves more than an increase in the NRA from 65 to 66.

\paragraph{The welfare cost of a dollar saved.} Once the reforms with infinite MVPF are exhausted, further savings require reforms that lower welfare, and the MVPF of such a reform measures its welfare cost per dollar saved. Raising the NRA costs 16 to 18 cents of willingness to pay, in dollars of consumption from age 62, per dollar of government spending saved. Our grid contains two other families against which this cost can be compared, Medicare coverage after 65 and income taxes after 62, but we simulate them as expansions, with MVPFs of 0.146 to 0.155 and 0.23 to 0.26. To a first order, the MVPF of a cut equals the MVPF of the corresponding expansion. For a small change $d\theta$ in a policy parameter $\theta$, both the WTP and the net cost are proportional to $d\theta$; reversing the sign of $d\theta$ reverses the sign of both and leaves their ratio unchanged. Read this way, raising the NRA costs about as much welfare per dollar saved as forgoing or reversing an expansion of Medicare coverage after 65, and less than raising income taxes on older individuals, which would cost about 23 to 26 cents per dollar raised. Among the welfare-reducing instruments we consider, NRA increases are thus a comparatively inexpensive source of revenue, even though they are dominated by rewards for delay for moderate savings.

This symmetry is a first-order argument, and our reforms are not marginal. If the WTP is concave in the policy parameter, as diminishing returns imply, the welfare lost by moving the parameter down exceeds the welfare gained by moving it up by the same amount, and the reversed MVPFs understate the welfare cost of a cut. For old-age income taxes, the MVPFs of I1--I5 hardly change with the size of the reform, which suggests that this curvature is mild over the range we consider. For Medicare, the gap between our estimates and the estimates of the MVPF of its introduction, 0.52 to 3.83 (Section~\ref{sec:results_medicare}), suggests that the curvature is pronounced over larger changes, so that cuts in coverage could cost more per dollar than our expansions return. Behavioral responses to cuts and expansions also need not be symmetric when participation involves fixed costs of work and a cost of reentry. We have not simulated the reverse reforms, and the comparison between NRA increases and tax increases or Medicare cuts is an approximation rather than a computed result. The comparison between NRA increases and reward-delay reforms involves no such approximation.

\paragraph{Who bears the gap.} Our analysis concerns a single cohort, facing reforms announced before age 30 and in force over its whole life. A financing gap also raises the question of how a legacy debt is shared across cohorts, and a cohort analysis does not determine who should bear it. The MVPF ranks the instruments for raising a dollar from a given cohort; it does not say which cohorts should contribute that dollar. The transition generations, cohorts already in the middle of their careers or near retirement when a reform is announced, are not modeled. For them, both the welfare cost and the fiscal saving of a reform depend on how much of their working life remains when it is announced: an NRA increase announced late in a career leaves little room to adjust saving and labor supply, and a higher credit can change the behavior only of those who have not yet claimed. Prices are held fixed throughout, so the analysis also abstracts from the effects of a consolidation on wages and interest rates (Section~\ref{sec:discussion}).

With these qualifications, the MVPFs suggest an order for closing a financing gap. Reforms that reward later claiming come first, because they save money and raise welfare. Beyond the savings they can deliver, the choice among reforms that lower welfare should follow their welfare cost per dollar saved, on which increases in the NRA compare favorably with taxes on older individuals.
